\chapter{自然数、归纳与递归}
\begin{qlnote}{阅读参照}Enderton, \emph{Set Theory}, Chapter 4.\end{qlnote}

\section{从无穷公理得到 \texorpdfstring{\(\omega\)}{omega}}
定义后继 \(S(x)=x\cup\{x\}\)。无穷公理给出某个归纳集合 \(I\)。
在 \(\mathcal P(I)\) 内以分离选出所有归纳子集，得到非空集合
\(\mathcal J=\{J\subseteq I:J\text{ 是归纳集合}\}\)，因为 \(I\in\mathcal J\)。
定义 \(\omega=\bigcap\mathcal J\)。

\begin{theorem}{最小归纳集合}
\(\omega\) 是归纳集合，而且包含于每个归纳集合。
\end{theorem}
\begin{proof}
每个 \(J\in\mathcal J\) 都含空集，所以交集含空集；后继闭合性同样保留于交集。
若 \(K\) 是任意归纳集合，则 \(I\cap K\) 是归纳子集，属于 \(\mathcal J\)。
故 \(\omega\subseteq I\cap K\subseteq K\)。
\end{proof}
这样避免把「所有归纳集合」未经证明地当成一个集合。
在这个构造中
\[
 0=\varnothing,\quad 1=\{0\},\quad 2=\{0,1\},\quad 3=\{0,1,2\},\quad n+1=S(n).
\]
每个自然数是其所有严格较小自然数组成的集合；这项性质须由归纳建立。

\section{归纳原理}
\begin{theorem}{Induction on \texorpdfstring{\(\omega\)}{omega}}
若 \(A\subseteq\omega\)，且 \(0\in A\)、\(n\in A\Rightarrow S(n)\in A\)，则 \(A=\omega\)。
\end{theorem}
\begin{proof}
条件说明 \(A\) 是归纳集合，由最小性有 \(\omega\subseteq A\)，结合反向包含得到相等。
\end{proof}
对公式 \(\varphi(n,\bar p)\)，先以分离得到满足该公式的自然数集合，再应用上述定理。
因此「对每个公式成立的归纳」是一族实例；它不是在一阶语言中直接量化于公式。

\section{递归定理与运算}
\begin{theorem}{Recursion on \texorpdfstring{\(\omega\)}{omega}}
给定集合 \(A\)、\(a\in A\) 及函数 \(g:A\to A\)，存在唯一函数 \(h:\omega\to A\)，满足
\(h(0)=a\) 和 \(h(S(n))=g(h(n))\)。
\end{theorem}
\begin{qlnote}{证明路线}
先对每个有限初始段归纳构造递归图像，并证明它们在共同定义域上一致。
这些图像是 \(\omega\times A\) 的子集；可在其幂集中分离出全部有限近似图像，
取并得到 \(h\)。归纳保证全域性，兼容性保证单值性；再用归纳证明两个满足条件的函数相同。
逐项补足有限初始段和近似图像的定义，是本章的原书阅读任务。
\end{qlnote}

由递归可定义
\[
 m+0=m,\quad m+S(n)=S(m+n),\qquad
 m\cdot0=0,\quad m\cdot S(n)=m\cdot n+m.
\]
定义之后还须用归纳证明结合律、交换律及相应消去律。
\textbf{归纳}建立性质；\textbf{递归}构造对象。递推式需要存在唯一性定理才能成为有效定义。
